\section{Residue-parity localization in the balancing-Pell packet}
\label{sec:pell-residue-parity-localization}

Write
\[
 (1+\sqrt2)^n=A_n+B_n\sqrt2 .
\]
For odd $n$ these coprime coordinates satisfy
$A_n^2-2B_n^2=-1$.  The preceding four-prime analysis proves that, at an
odd prime index $\ell$, the odd-exponent prime sets $O_A,O_B$ obey
\begin{equation}\label{eq:pell-cross-character-product-recalled}
 \prod_{q\in O_A}\prod_{r\in O_B}
       \left(\frac qr\right)=\left(\frac2\ell\right).
\end{equation}
We now localize the depth-three primes by congruence class.

Multiplication by $1+\sqrt2$ gives
\[
 A_{n+1}=A_n+2B_n,\qquad B_{n+1}=A_n+B_n.
\]
Since
\[
 (1+\sqrt2)^8=577+408\sqrt2\equiv1\pmod8,
\]
the coordinate pair is periodic modulo eight.  Direct evaluation gives
\begin{equation}\label{eq:pell-mod-eight-table}
\begin{array}{c|cccc}
n\pmod8&1&3&5&7\\ \hline
A_n\pmod8&1&7&1&7\\
B_n\pmod8&1&5&5&1.
\end{array}
\end{equation}
If an odd prime $q$ divides $A_n$, the norm equation gives
$2B_n^2\equiv1\pmod q$, whence $q\equiv1,7\pmod8$.  If an odd prime $r$
divides $B_n$, then $A_n^2\equiv-1\pmod r$, whence
$r\equiv1,5\pmod8$.

We use the following finite-product observation.  If $\rho^2=1$ and
$\rho\ne1$ in a commutative monoid, every $x_i$ is $1$ or $\rho$, and
\[
                       \prod_i x_i^{e_i}=\rho,
\]
then some $x_i=\rho$ has odd exponent.  Otherwise each nontrivial factor
would occur to even exponent and the product would be one.  If every
$e_i\ge2$, the selected exponent is at least three.

\begin{theorem}[Forced-residue depth-three localization]
\label{thm:pell-forced-residue-depth-three}
Let $\ell$ be an odd prime and suppose $A_\ell B_\ell$ is squarefull.
Then:
\begin{enumerate}
\item if $\ell\equiv3,7\pmod8$, some prime $q_A\mid A_\ell$ occurs to an
odd exponent at least three and satisfies $q_A\equiv7\pmod8$;
\item if $\ell\equiv3,5\pmod8$, some prime $q_B\mid B_\ell$ occurs to an
odd exponent at least three and satisfies $q_B\equiv5\pmod8$.
\end{enumerate}
\end{theorem}

\begin{proof}
The residue table shows that both coordinates are odd.  Reduce the two
prime-power factorizations modulo eight.  The allowed local
classes are $\{1,7\}$ and $\{1,5\}$, respectively.  In the stated index
classes, \eqref{eq:pell-mod-eight-table} makes the total product the
nontrivial involution $7$ or $5$.  The finite-product observation selects an
odd exponent.  Since $\gcd(A_\ell,B_\ell)=1$, squarefullness of
$A_\ell B_\ell$ makes every positive exponent in either channel at least
two, so the selected odd exponent is at least three.
\end{proof}

For $\ell\equiv3\pmod8$, both conclusions hold simultaneously.  Thus the
two previously unspecified deep primes can be required to occupy the
$7$-class in the $A$ channel and the $5$-class in the $B$ channel.

The Lean statement uses an explicit factor-list packet interface.  Each
entry stores a prime and its exponent; the packet records the allowed local
classes and the squarefull lower bound.  If the $A$ residue product is $7$
or the $B$ residue product is $5$, the corresponding list contains the
prime and odd exponent asserted above.  Thus the formal theorem combines
the residue-product hypothesis and the depth conclusion in one statement,
while passage from an actual Pell factorization to that interface is the
elementary argument in this proof.

\begin{theorem}[A forced cross-channel nonresidue]
\label{thm:pell-forced-cross-nonresidue}
If $\ell\equiv3,5\pmod8$ and $A_\ell B_\ell$ is squarefull, then there are
odd-exponent primes $q\mid A_\ell$ and $r\mid B_\ell$, each of exponent at
least three, such that
\[
                            \left(\frac qr\right)=-1.
\]
\end{theorem}

\begin{proof}
Here $(2/\ell)=-1$.  Equation
\eqref{eq:pell-cross-character-product-recalled} is a finite product of
nonzero signs with value $-1$: coprimality of the two channels prevents a
zero Legendre symbol.  Hence at least one factor is negative.  Membership in
$O_A\times O_B$ gives odd exponents; squarefullness makes them at least
three.
\end{proof}

Here too the Lean statement is combined rather than left as two unrelated
lemmas: a flattened factor-list packet stores every pair in
$O_A\times O_B$, its two odd squarefull exponents, and its actual Legendre
symbol.  A negative product yields a stored pair with symbol $-1$ and both
exponents at least three.

The existential quantifier is sharp under the ordinary actual prime-index
premises.  At $\ell=11$ one has
\[
 A_{11}=8119=23\cdot353,\qquad B_{11}=5741,
\]
with all three factors prime, while
\[
 \left(\frac{23}{5741}\right)=-1,
 \qquad
 \left(\frac{353}{5741}\right)=1.
\]
Indeed, $252^2\equiv353\pmod {5741}$, whereas quadratic reciprocity reduces
the first symbol to $\left(\frac{14}{23}\right)=-1$.  The product is still
$-1=\left(\frac2{11}\right)$ because $11\equiv3\pmod8$.  The factorization shows that
$23,353\in O_A$ and $5741\in O_B$.  This is a full actual counterexample to
the strengthening that every cross pair is a nonresidue.  It has three exponent-one factors,
so it does not satisfy the squarefull premise and does not close the Pell
route.

The companion Lean module checks the integral recurrence and norm, the
eight-step formula, all residue classes in
\eqref{eq:pell-mod-eight-table}, the abstract odd-exponent and depth-three
extraction, the two explicit factor-list packet theorems, and one combined
index-eleven certificate containing the prime-index premise, global
character, coordinates, factorization, exponent-one witnesses and Legendre
values.  The open point remains an actual
opposite-channel depth-three exclusion; the new residue and character
conditions narrow that packet without yet making it inconsistent.
